\begin{ouverture}[grandeurs]{11}{Longueurs, masses, contenances}{Partie 3 · Grandeurs et mesures}
{Conversions raisonnées · Estimations · Problèmes}
\objectif{Convertir en s’appuyant sur les relations entre unités.}
\objectif{Estimer des mesures avec des références fiables.}
\objectif{Résoudre des problèmes de mesures en plusieurs étapes.}
\prerequis{Unités du CM1 et décimaux jusqu’aux millièmes.}
\end{ouverture}
\pagestyle{otMain}\essentiel\fontsize{11}{15}\selectfont
\phantomsection\label{cours-11-1}
\essSection{1}{Les relations entre unités}
\begin{propriete}Longueurs : $1\ \mathrm{km}=1\,000\ \mathrm m$, $1\ \mathrm m=100\ \mathrm{cm}=1\,000\ \mathrm{mm}$. Masses : $1\ \mathrm t=1\,000\ \mathrm{kg}$, $1\ \mathrm{kg}=1\,000\ \mathrm g$, $1\ \mathrm g=1\,000\ \mathrm{mg}$. Contenances : $1\ \mathrm{hL}=100\ \mathrm L$, $1\ \mathrm L=100\ \mathrm{cL}=1\,000\ \mathrm{mL}$.\end{propriete}
\begin{savoirfaire}[Convertir sans tableau]Convertir, c’est citer la relation : $3{,}5$ m $=350$ cm car 1 m vaut 100 cm ; $0{,}75$ kg $=750$ g car 1 kg vaut $1\,000$ g ; $2\,500$ mL $=2{,}5$ L. Les millièmes rendent les conversions fines : 1 mm $=0{,}001$ m.\end{savoirfaire}
\begin{exemple}Références utiles : une porte 2 m, un grand pas 1 m, une pomme 150 g, une pièce de 1 euro $7{,}5$ g, une canette 33 cL, une baignoire 150 L, un litre d’eau pèse environ 1 kg.\end{exemple}
\begin{attention}Comparer exige la même unité : $1{,}2$ km contre 980 m se joue en mètres. Une réponse s’écrit avec son unité, choisie pour rester lisible : 1 250 m se dit aussi $1{,}25$ km.\end{attention}
\clearpage
\phantomsection\label{cours-11-2}
\essSection{2}{Des mesures qui se combinent}
\begin{savoirfaire}[Même unité avant tout]Additionner des mesures : tout mettre dans la même unité d’abord. $1{,}2$ kg $+850$ g $=1\,200+850=2\,050$ g $=2{,}05$ kg.\end{savoirfaire}
\begin{exemple}Un périmètre se calcule dans une unité unique : un champ de 125 m sur $87{,}5$ m a pour tour $2\times(125+87{,}5)=425$ m.\end{exemple}
\begin{propriete}Les mesures se partagent : $1{,}5$ L réparti dans 6 verres donne $150:6=25$ cL par verre. Les relations entre unités transforment souvent une division décimale en division entière.\end{propriete}
\begin{attention}Estimer avant de répondre : une baleine de 30 g, un cartable de 300 kg ou un verre de 20 L sont invraisemblables. La vraisemblance fait partie de la réponse.\end{attention}

\clearpage\rubrique{Découverte}
\begin{blocLivre}{ex-11-1}\begin{exercice}[1]{Relations à citer}
Convertir en citant la relation utilisée : $3{,}5$ m en cm ; $0{,}75$ kg en g ; $2\,500$ mL en L.
\end{exercice}\end{blocLivre}
\begin{blocLivre}{ex-11-2}\begin{exercice}[1]{Vers les millièmes}
Convertir : 45 mm en m ; 3 g en kg ; 7 mL en L.
\end{exercice}\end{blocLivre}
\begin{blocLivre}{ex-11-3}\begin{exercice}[1]{Comparer}
Comparer : $1{,}2$ km et 980 m ; $0{,}5$ t et 520 kg ; 75 cL et $0{,}7$ L.
\end{exercice}\end{blocLivre}
\begin{blocLivre}{ex-11-4}\begin{exercice}[1]{Estimations}
Choisir la mesure vraisemblable : masse d’un vélo (12 kg ou 120 kg ?) ; contenance d’une tasse (25 cL ou 25 L ?) ; longueur d’une salle de classe (8 m ou 80 m ?).
\end{exercice}\end{blocLivre}
\begin{blocLivre}{ex-11-5}\begin{exercice}[1]{Additionner des mesures}
Calculer : $1{,}2$ kg $+$ 850 g, en grammes puis en kilogrammes.
\end{exercice}\end{blocLivre}

\rubrique{Entraînement}
\begin{blocLivre}{ex-11-6}\begin{exercice}[2]{Le tour du champ}
Un champ rectangulaire mesure 125 m sur $87{,}5$ m. Calculer son périmètre, en mètres puis en kilomètres.
\end{exercice}\end{blocLivre}
\begin{blocLivre}{ex-11-7}\begin{exercice}[2]{La carafe et les verres}
Une carafe contient $1{,}5$ L. On remplit 6 verres identiques en la vidant exactement. Quelle est la contenance d’un verre, en centilitres ?
\end{exercice}\end{blocLivre}
\begin{blocLivre}{ex-11-8}\begin{exercice}[2]{Le colis à la poste}
Un colis contient 3 livres de 650 g chacun et un emballage de 120 g. Peut-on l’envoyer au tarif « moins de 2 kg » ?
\end{exercice}\end{blocLivre}
\begin{blocLivre}{ex-11-9}\begin{exercice}[2]{L’eau de la semaine}
Une famille boit 2 packs de 6 bouteilles de $1{,}5$ L par semaine. Quel volume d’eau cela fait-il ? Combien de litres par jour, en moyenne régulière sur 7 jours ? Donner une valeur approchée au dixième.
\end{exercice}\end{blocLivre}
\begin{blocLivre}{ex-11-10}\begin{exercice}[2]{Le médicament}
Un flacon de sirop contient 200 mL. La dose est de 5 mL, trois fois par jour. Combien de jours complets le flacon dure-t-il ?
\end{exercice}\end{blocLivre}

\rubrique{Problèmes}
\begin{blocLivre}{ex-11-11}\begin{exercice}[3]{La citerne du jardin}
Une citerne de 3 hL est pleine. On arrose avec 12 arrosoirs de 9 L, puis on remplit un bassin de 150 L. Quelle contenance d’eau reste-t-il ?
\end{exercice}\end{blocLivre}
\begin{blocLivre}{ex-11-12}\begin{exercice}[3]{Les sacs de farine}
Un boulanger reçoit $0{,}5$ t de farine en sacs de 25 kg. Combien de sacs reçoit-il ?
\end{exercice}\end{blocLivre}
\begin{blocLivre}{ex-11-13}\begin{exercice}[3]{Le ruban exact}
Pour un paquet, il faut 65 cm de ruban. Combien de paquets complets prépare-t-on avec une bobine de 5 m ? Quelle longueur reste-t-il ?
\end{exercice}\end{blocLivre}
\begin{blocLivre}{ex-11-14}\begin{exercice}[3]{La tonne en question}
Un ascenseur de chantier accepte $0{,}8$ t. On veut monter 15 sacs de ciment de 35 kg et 4 seaux de 12 kg chacun. Est-ce possible en un seul voyage ?
\end{exercice}\end{blocLivre}
\begin{blocLivre}{ex-11-15}\begin{exercice}[3]{L’étiquette douteuse}
Une étiquette annonce : « Bouteille de 75 cL, soit $7{,}5$ L ». Corriger l’erreur et donner la bonne équivalence en litres et en millilitres.
\end{exercice}\end{blocLivre}

\clearpage
